\documentclass[12pt]{ctexart}
\usepackage{amsmath,amssymb}
\usepackage{graphicx}
\usepackage{geometry}
\usepackage{xcolor}
\usepackage{booktabs}
\usepackage{longtable}
\usepackage{array}
\usepackage{enumitem}
\usepackage{float}
\usepackage{caption}
\usepackage{tikz}
\usetikzlibrary{positioning,arrows.meta,calc,shapes.misc,fit,backgrounds}
\usepackage{tcolorbox}
\usepackage{titlesec}
\usepackage{fancyhdr}
\usepackage{setspace}
\usepackage{microtype}
\usepackage{hyperref}
\hypersetup{
  colorlinks=true,
  linkcolor=blue,
  citecolor=darkgray,
  urlcolor=blue,
  bookmarks=true,
  bookmarksnumbered=true,
  unicode=true,
  pdftitle={拉普拉斯变换、收敛域与传递函数：控制论入门讲义（附两轮平衡车控制框图）},
  pdfsubject={控制论基础讲义},
  pdfkeywords={拉普拉斯变换, 传递函数, 方框图, 两轮平衡车, PID}
}
\setmainfont{Noto Serif}
\setsansfont{Noto Sans}
\setmonofont{DejaVu Sans Mono}
\setCJKmainfont{Noto Serif CJK SC}
\setCJKsansfont{Noto Sans CJK SC}
\setCJKmonofont{Noto Sans Mono CJK SC}
\geometry{a4paper, top=2.35cm, bottom=2.35cm, left=2.55cm, right=2.35cm}
\setstretch{1.16}
\setlength{\parindent}{2em}
\setlength{\parskip}{0.35em}
\setlist[itemize]{itemsep=0.25em, leftmargin=1.8em}
\setlist[enumerate]{itemsep=0.25em, leftmargin=1.8em}
\captionsetup{font=small, labelfont=bf}
\titleformat{\section}{\Large\bfseries}{\thesection}{0.8em}{}
\titleformat{\subsection}{\large\bfseries}{\thesubsection}{0.8em}{}
\titlespacing*{\section}{0pt}{1.2em}{0.7em}
\titlespacing*{\subsection}{0pt}{1.0em}{0.5em}
\providecommand{\tightlist}{\setlength{\itemsep}{0pt}\setlength{\parskip}{0pt}}
\setlength{\headheight}{14pt}
\setcounter{secnumdepth}{-2}
\pagestyle{fancy}
\fancyhf{}
\fancyhead[L]{\small 拉普拉斯变换与传递函数}
\fancyhead[R]{\small 附两轮平衡车控制框图}
\fancyfoot[C]{\small \thepage}
\renewcommand{\headrulewidth}{0.4pt}

\title{拉普拉斯变换、收敛域与传递函数：\\控制论入门讲义（附两轮平衡车控制框图）}
\author{}
\date{}
\begin{document}
\maketitle
\tableofcontents
\newpage
\hypertarget{ux4e00ux62c9ux666eux62c9ux65afux53d8ux6362ux7684ux5b9aux4e49ux4e0eux5b58ux5728ux6761ux4ef6}{%
\subsection{一、拉普拉斯变换的定义与存在条件}\label{ux4e00ux62c9ux666eux62c9ux65afux53d8ux6362ux7684ux5b9aux4e49ux4e0eux5b58ux5728ux6761ux4ef6}}

\hypertarget{ux4e24ux79cdux5b9aux4e49}{%
\subsubsection{1.1 两种定义}\label{ux4e24ux79cdux5b9aux4e49}}

控制论中使用单边拉普拉斯变换：

\[
F(s)=\int_{0^{-}}^{\infty} f(t)\,e^{-st}\,dt
\]

其中 \(s=\sigma+j\omega\) 是复变量。积分下限取 \(0^{-}\) 而不是
\(0\)，是为了把 \(t=0\)
处可能出现的冲激（例如阶跃输入的导数）完整包含进积分。控制论研究的是从某一初始时刻起外加输入、观察系统响应的问题，时间轴天然地从
\(0\) 开始，因此单边变换是默认工具。

信号与系统课程中使用双边拉普拉斯变换：

\[
F(s)=\int_{-\infty}^{\infty} f(t)\,e^{-st}\,dt
\]

它处理的是整条时间轴上的信号，包括 \(t<0\)
的部分。本文后面讨论收敛域时会反复用到双边变换，因为只有在双边框架下，收敛域的必要性才暴露得最清楚。

\hypertarget{ux5b58ux5728ux6761ux4ef6ux6307ux6570ux9636}{%
\subsubsection{1.2
存在条件：指数阶}\label{ux5b58ux5728ux6761ux4ef6ux6307ux6570ux9636}}

积分 \(\int f(t)e^{-st}dt\) 不是对任何 \(f(t)\)
都有意义。使变换存在的一个充分条件是 \(f(t)\) 为指数阶函数：存在常数
\(M>0\) 与实数 \(a\)，使得 \(t\) 充分大时

\[
|f(t)|\le Me^{at}.
\]

严格地说，这个充分条件还有一半：\(f(t)\)
应当分段连续，即在任意有限区间上只有有限个间断点，且每个间断点处的左右极限都存在。分段连续保证积分在任何有限区间上有意义，指数阶保证积分在无穷远处收敛，两半合起来变换才存在。

多项式、指数函数、正弦函数都是指数阶的。\(e^{t^{2}}\)
不是指数阶的：对任何固定的 \(a\)，当 \(t\) 充分大时
\(e^{t^{2}}>Me^{at}\)，因此 \(e^{t^{2}}\) 的拉普拉斯变换不存在。

\hypertarget{ux6536ux655bux6761ux4ef6ux7684ux63a8ux5bfc}{%
\subsubsection{1.3
收敛条件的推导}\label{ux6536ux655bux6761ux4ef6ux7684ux63a8ux5bfc}}

设 \(f(t)\) 为指数阶，\(|f(t)|\le Me^{at}\)。把 \(s=\sigma+j\omega\)
代入被积函数取模：

\[
\left|e^{-st}f(t)\right|=\left|e^{-\sigma t}\right|\cdot\left|e^{-j\omega t}\right|\cdot|f(t)|=e^{-\sigma t}|f(t)|\le Me^{-(\sigma-a)t}.
\]

这里用到 \(\left|e^{-j\omega t}\right|=1\)：虚部 \(j\omega\)
只贡献一个模恒为 \(1\) 的旋转因子，对积分是否收敛毫无影响。当
\(\sigma>a\) 时，\(e^{-(\sigma-a)t}\to 0\)（\(t\to\infty\)），积分
\(\int_{0}^{\infty}Me^{-(\sigma-a)t}dt=\dfrac{M}{\sigma-a}\)
有限，由比较判别法，原积分绝对收敛。

这个推导给出两个结论。第一，变换存在的条件是
\(\mathrm{Re}(s)=\sigma>a\)，即收敛与否只取决于 \(s\)
的实部。第二，收敛域必然是 \(s\)
平面上一条竖直线右侧的区域（对双边变换则可能是带状或左侧区域），因为判据里根本不出现
\(\omega\)。收敛域是竖直区域这一点，后面会反复用到。

\hypertarget{ux540cux4e00ux6761ux4ef6ux4fddux8bc1ux5faeux5206ux6027ux8d28ux6210ux7acb}{%
\subsubsection{1.4
同一条件保证微分性质成立}\label{ux540cux4e00ux6761ux4ef6ux4fddux8bc1ux5faeux5206ux6027ux8d28ux6210ux7acb}}

微分性质 \(\mathcal{L}\{f'(t)\}=sF(s)-f(0^{-})\) 由分部积分得到：

\[
\int_{0^{-}}^{\infty}f'(t)e^{-st}dt=\left.e^{-st}f(t)\right|_{0^{-}}^{\infty}+s\int_{0^{-}}^{\infty}f(t)e^{-st}dt.
\]

边界项在上限处为 \(\lim_{t\to\infty}e^{-st}f(t)\)。由 1.3
的估计，\(\left|e^{-st}f(t)\right|\le Me^{-(\sigma-a)t}\)，当
\(\sigma>a\) 时该边界项等于 \(0\)；下限处代入得 \(-f(0^{-})\)。于是

\[
\mathcal{L}\{f'(t)\}=sF(s)-f(0^{-}).
\]

注意使变换收敛的条件与使边界项消失的条件是同一个条件
\(\sigma>a\)。这不是巧合：分部积分的合法性本来就依赖被积函数在无穷远处衰减。后面会看到，``收敛域''这个概念所做的全部工作，就是把这类条件从隐含变成显式。

\hypertarget{ux4e8cux6536ux655bux57dfux5b83ux4eceux54eaux91ccux6765ux7ba1ux4ec0ux4e48}{%
\subsection{二、收敛域：它从哪里来，管什么}\label{ux4e8cux6536ux655bux57dfux5b83ux4eceux54eaux91ccux6765ux7ba1ux4ec0ux4e48}}

\hypertarget{ux5b8cux6574ux7684ux53d8ux6362ux662fux5f0fux5b50ux52a0ux6536ux655bux57df}{%
\subsubsection{2.1
完整的变换是''式子加收敛域''}\label{ux5b8cux6574ux7684ux53d8ux6362ux662fux5f0fux5b50ux52a0ux6536ux655bux57df}}

对指数信号 \(e^{at}u(t)\)（\(u(t)\) 为单位阶跃），其双边拉普拉斯变换为

\[
\int_{0}^{\infty}e^{at}e^{-st}dt=\int_{0}^{\infty}e^{-(s-a)t}dt=\frac{1}{s-a},\qquad \mathrm{Re}(s)>a.
\]

收敛条件来自 \(e^{-(s-a)t}\) 在 \(t\to\infty\) 时衰减，即
\(\mathrm{Re}(s-a)>0\)。

再看左边信号 \(-e^{at}u(-t)\)，它在 \(t<0\) 时等于
\(-e^{at}\)，\(t\ge 0\) 时为 \(0\)：

\[
\int_{-\infty}^{0}-e^{at}e^{-st}dt=-\int_{-\infty}^{0}e^{-(s-a)t}dt=\frac{1}{s-a},\qquad \mathrm{Re}(s)<a.
\]

这次积分区间是 \((-\infty,0]\)，收敛要求 \(e^{-(s-a)t}\) 在
\(t\to-\infty\) 时衰减，即 \(\mathrm{Re}(s-a)<0\)。

两个完全不同的信号，变换的式子都是
\(\dfrac{1}{s-a}\)，唯一的区别是收敛域：前者为
\(\mathrm{Re}(s)>a\)，后者为
\(\mathrm{Re}(s)<a\)。这说明拉普拉斯变换的完整定义是''式子加收敛域''，二者缺一不可。只给式子不给收敛域，反变换不唯一。

\hypertarget{ux7b2cux4e00ux4e2aux8befux533aux6536ux655bux57dfux9650ux5236ux7684ux662f-sux4e0dux662f-t}{%
\subsubsection{\texorpdfstring{2.2 第一个误区：收敛域限制的是
\(s\)，不是
\(t\)}{2.2 第一个误区：收敛域限制的是 s，不是 t}}\label{ux7b2cux4e00ux4e2aux8befux533aux6536ux655bux57dfux9650ux5236ux7684ux662f-sux4e0dux662f-t}}

初学者常问：``超出作用域的地方怎么处理？''这个问法默认收敛域是对时间
\(t\) 的限制，仿佛 \(t\)
跑出某个范围后解就失效了。事实不是这样。收敛域是 \(s\)
平面上的区域，它记录的是''积分在哪些 \(s\) 上有意义''，与 \(t\)
的取值无关。反变换

\[
f(t)=\frac{1}{2\pi j}\int_{\sigma-j\infty}^{\sigma+j\infty}F(s)e^{st}ds
\]

沿着收敛域内一条竖直线积分，得到的结果对全部 \(t\)（双边情形为
\(-\infty<t<\infty\)）有效。没有任何一段解被''扔掉''。收敛域限制的是
\(s\) 平面，不是时间轴。

\hypertarget{ux7b2cux4e8cux4e2aux8befux533aux6536ux655bux57dfux4e0dux662fux9644ux52a0ux7684ux6280ux672fux7ec6ux8282}{%
\subsubsection{2.3
第二个误区：收敛域不是附加的技术细节}\label{ux7b2cux4e8cux4e2aux8befux533aux6536ux655bux57dfux4e0dux662fux9644ux52a0ux7684ux6280ux672fux7ec6ux8282}}

既然只给式子反变换不唯一，收敛域就不是可以随手丢掉的技术细节，而是变换本身的一半。2.1
的例子表明，同一个 \(\dfrac{1}{s-a}\)
配上右侧收敛域是因果指数，配上左侧收敛域是反因果指数。式子相同而收敛域不同，对应的是物理上截然不同的信号。

\hypertarget{ux6536ux655bux57dfux7684ux56dbux6761ux89c4ux5219}{%
\subsubsection{2.4
收敛域的四条规则}\label{ux6536ux655bux57dfux7684ux56dbux6761ux89c4ux5219}}

以下是双边变换下有理性变换（\(s\) 的有理分式）的收敛域规则，均可由 1.3
的推导直接得到：

\begin{enumerate}
\def\labelenumi{\arabic{enumi}.}
\tightlist
\item
  收敛域内不含极点，收敛域的边界由极点围成。理由在收敛条件的形状：收敛条件形如
  \(\mathrm{Re}(s)>a\) 或 \(\mathrm{Re}(s)<a\)，其中 \(a\)
  恰为某个极点的实部，因此收敛域的边界只能是过极点的竖直线；收敛域本身则是这些竖直线划分出的竖直区域------最右极点以右、最左极点以左，或相邻两极点之间的带状区域------无论哪种情形都不含任何极点。
\item
  和式的收敛域至少包含各项收敛域的交集。和的积分要收敛，只需各项积分同时收敛，因此各项收敛域的交集一定落在和式的收敛域内。但和式的收敛域可能严格更大：相加时若发生零极点相消，某个极点会消失，由它划定的边界随之解除。例如
  \(X_1(s)=\dfrac{1}{s+1}\)（收敛域 \(\mathrm{Re}(s)>-1\)）与
  \(X_2(s)=-\dfrac{1}{s+1}+\dfrac{1}{s+2}\)（收敛域
  \(\mathrm{Re}(s)>-1\)）相加，得 \(X_1(s)+X_2(s)=\dfrac{1}{s+2}\)，极点
  \(-1\) 被消去，收敛域从交集 \(\mathrm{Re}(s)>-1\) 扩大为
  \(\mathrm{Re}(s)>-2\)。
\item
  逐项反变换的判定：对某一项 \(\dfrac{1}{s-a}\)，若总收敛域位于该项极点
  \(a\) 的右侧，该项反变换为 \(e^{at}u(t)\)；若位于左侧，该项反变换为
  \(-e^{at}u(-t)\)。这正是 2.1 两个积分的结论。
\item
  附加条件与收敛域的对应：因果对应最右极点以右的收敛域；稳定对应收敛域包含虚轴；反因果对应最左极点以左的收敛域。这些条件可能互相矛盾：若系统有右半平面极点，则''因果''要求收敛域在该极点右侧，``稳定''要求收敛域含虚轴，二者不能同时成立，``因果且稳定''无解。
\end{enumerate}

规则 4 中''稳定对应含虚轴''的严格理由在 3.3 节给出，此处先作为规则列出。

\hypertarget{ux5173ux952eux793aux4f8bhsdfrac1s1s-2-ux7684ux4e09ux79cdux6536ux655bux57df}{%
\subsubsection{\texorpdfstring{2.5
关键示例：\(H(s)=\dfrac{1}{(s+1)(s-2)}\)
的三种收敛域}{2.5 关键示例：H(s)=\textbackslash dfrac\{1\}\{(s+1)(s-2)\} 的三种收敛域}}\label{ux5173ux952eux793aux4f8bhsdfrac1s1s-2-ux7684ux4e09ux79cdux6536ux655bux57df}}

先做部分分式展开。设

\[
\frac{1}{(s+1)(s-2)}=\frac{A}{s-2}+\frac{B}{s+1}.
\]

两边同乘 \((s+1)(s-2)\)，得 \(1=A(s+1)+B(s-2)\)。比较两边系数：\(s\)
的系数给出 \(A+B=0\)，常数项给出 \(A-2B=1\)，解得
\(A=\dfrac{1}{3}\)，\(B=-\dfrac{1}{3}\)。同一结果也可用留数法直接读出：\(A=\lim_{s\to 2}(s-2)H(s)=\dfrac{1}{2+1}=\dfrac{1}{3}\)，\(B=\lim_{s\to -1}(s+1)H(s)=\dfrac{1}{-1-2}=-\dfrac{1}{3}\)。于是

\[
H(s)=\frac{1}{(s+1)(s-2)}=\frac{1}{3}\left(\frac{1}{s-2}-\frac{1}{s+1}\right).
\]

极点为 \(s=2\) 与 \(s=-1\)。两条竖直线 \(\mathrm{Re}(s)=2\) 与
\(\mathrm{Re}(s)=-1\) 把 \(s\)
平面分成三个竖直区域，每个区域都是一种合法的收敛域。

\textbf{情形一：\(\mathrm{Re}(s)>2\)。}
收敛域在最右极点以右，对应因果系统。对每一项分别判定：收敛域在极点 \(2\)
的右侧，\(\dfrac{1}{s-2}\) 反变换为 \(e^{2t}u(t)\)；收敛域在极点 \(-1\)
的右侧，\(\dfrac{1}{s+1}\) 反变换为 \(e^{-t}u(t)\)。合起来：

\[
h(t)=\frac{1}{3}\left(e^{2t}-e^{-t}\right)u(t).
\]

含 \(e^{2t}\)，\(t\to\infty\) 时发散，系统不稳定。

\textbf{情形二：\(-1<\mathrm{Re}(s)<2\)。}
收敛域是两条竖直线之间的带状区域。逐项判定：收敛域在极点 \(2\)
的左侧，\(\dfrac{1}{s-2}\) 反变换为 \(-e^{2t}u(-t)\)；收敛域在极点
\(-1\) 的右侧，\(\dfrac{1}{s+1}\) 反变换为 \(e^{-t}u(t)\)。合起来：

\[
h(t)=-\frac{1}{3}e^{-t}u(t)-\frac{1}{3}e^{2t}u(-t).
\]

这是双边信号：\(t\ge 0\) 时只有 \(-\frac{1}{3}e^{-t}\)，\(t<0\) 时只有
\(-\frac{1}{3}e^{2t}\)。\(t\to+\infty\) 时
\(e^{-t}\to 0\)，\(t\to-\infty\) 时
\(e^{2t}\to 0\)，冲激响应绝对可积，系统稳定。题目若声明''系统稳定''，必须取这个收敛域。

\textbf{情形三：\(\mathrm{Re}(s)<-1\)。}
收敛域在最左极点以左，对应反因果系统。逐项判定：收敛域在极点 \(2\)
的左侧，\(\dfrac{1}{s-2}\) 反变换为 \(-e^{2t}u(-t)\)；收敛域在极点
\(-1\) 的左侧，\(\dfrac{1}{s+1}\) 反变换为 \(-e^{-t}u(-t)\)。合起来：

\[
h(t)=\frac{1}{3}\left(e^{-t}-e^{2t}\right)u(-t).
\]

\(t\to-\infty\) 时 \(e^{-t}\) 发散，系统不稳定。

三种收敛域给出三个不同的
\(h(t)\)，不给附加条件时三种答案都合法。控制论题目默认因果，即取情形一。

\hypertarget{ux540cux4e00ux4e2a-sux4f5cux7528ux57dfux4e0dux8be5ux76f8ux540cux5417}{%
\subsubsection{\texorpdfstring{2.6 ``同一个
\(s\)，作用域不该相同吗''}{2.6 ``同一个 s，作用域不该相同吗''}}\label{ux540cux4e00ux4e2a-sux4f5cux7528ux57dfux4e0dux8be5ux76f8ux540cux5417}}

对上述逐项判定有一个自然的追问：整个 \(H(s)\)
的收敛域只有一个（例如情形二的带状区域），为什么反变换时每一项的''作用域''不一样？

澄清如下：整个 \(H(s)\)
确实只有一个收敛域，逐项判定时用的也是这同一个区域。判定规则不是给每一项各找一个收敛域，而是问同一个问题------这个区域位于该项极点的左侧还是右侧------对每一项分别回答。在情形二中，带状区域
\(-1<\mathrm{Re}(s)<2\) 对极点 \(2\) 而言在左侧，对极点 \(-1\)
而言在右侧，所以两项一项取左边信号、一项取右边信号。区域只有一个，回答有两个，因为极点有两个。

\hypertarget{ux4e3aux4ec0ux4e48ux6709ux4e9bux65f6ux5019ux505aux9898ux65f6ux53efux4ee5ux4e0dux7ba1ux6536ux655bux57df}{%
\subsubsection{2.7
为什么有些时候做题时可以不管收敛域}\label{ux4e3aux4ec0ux4e48ux6709ux4e9bux65f6ux5019ux505aux9898ux65f6ux53efux4ee5ux4e0dux7ba1ux6536ux655bux57df}}

原因有二。第一，控制论经常用单边变换，单边变换处理的信号都在 \(t\ge 0\)
上有定义，都是因果的。第二，由规则
4，因果信号的收敛域固定为''最右极点以右''。两个因素合起来，收敛域被唯一确定且自动满足，解题时写出来不增加任何信息，不写也不会错。

但''不用写''不等于''不存在''。一旦离开控制论的默认设定------处理双边信号、处理非因果系统、或者题目只给
\(H(s)\) 的式子而附加''稳定''条件------收敛域立刻变成答案的一部分。2.5
的情形二就是这种题目：只写 \(h(t)=\frac{1}{3}(e^{2t}-e^{-t})u(t)\)
是错的，因为那个响应不稳定。

\hypertarget{ux7279ux5f81ux6839ux6cd5ux4e3aux4ec0ux4e48ux6ca1ux6709ux6536ux655bux57df}{%
\subsubsection{2.8
特征根法为什么没有收敛域}\label{ux7279ux5f81ux6839ux6cd5ux4e3aux4ec0ux4e48ux6ca1ux6709ux6536ux655bux57df}}

特征根法解 \(\ddot y+a_1\dot y+a_0 y=0\) 时设
\(y=e^{rt}\)，代入得特征方程
\(r^2+a_1 r+a_0=0\)，求根后写出通解。全程没有积分，只有代数运算与求导，而代数运算与求导不存在收敛问题，所以特征根法里不会出现收敛域。

代价是：特征根法隐含地假设了解的因果性。它给出的是 \(t\ge 0\)
方向上的解（配合初始条件确定常数），它无法表示 \(-e^{at}u(-t)\)
这样的左边信号------这类信号根本不在''初始条件加右向演化''的框架之内。拉普拉斯变换（双边）把处理整条时间轴的信号作为目标，就必须为每个变换显式标明它描述的是哪一类信号，这个标注就是收敛域。换句话说，特征根法把因果性作为默认假设藏在方法内部，拉普拉斯变换把同一个假设显式化为收敛域。两种方法做的是同一件事，只是一个隐式、一个显式。

\hypertarget{ux4e09ux53d8ux6362ux7684ux672cux8d28ux6295ux5f71ux4e0eux7279ux5f81ux51fdux6570}{%
\subsection{三、变换的本质：投影与特征函数}\label{ux4e09ux53d8ux6362ux7684ux672cux8d28ux6295ux5f71ux4e0eux7279ux5f81ux51fdux6570}}

\hypertarget{ux5085ux91ccux53f6ux53d8ux6362ux5411ux6b63ux4ea4ux57faux6295ux5f71}{%
\subsubsection{3.1
傅里叶变换：向正交基投影}\label{ux5085ux91ccux53f6ux53d8ux6362ux5411ux6b63ux4ea4ux57faux6295ux5f71}}

傅里叶变换

\[
F(j\omega)=\int_{-\infty}^{\infty}f(t)e^{-j\omega t}dt
\]

可以理解为函数空间中的内积。把两个信号的内积定义为

\[
\langle x(t),\,y(t)\rangle=\int_{-\infty}^{\infty}x(t)\,y^{*}(t)\,dt,
\]

其中星号表示复共轭，则傅里叶变换正是信号与基函数 \(e^{j\omega t}\)
的内积：

\[
F(j\omega)=\langle f(t),\,e^{j\omega t}\rangle,
\]

即把信号 \(f(t)\) 向基函数 \(e^{j\omega t}\) 投影，投影系数
\(F(j\omega)\) 度量信号中频率为 \(\omega\) 的成分有多少。不同 \(\omega\)
的基函数彼此正交，因此傅里叶变换是信号在一组正交基下的坐标表示。它存在的充分条件是
\(f(t)\) 绝对可积，即 \(\int_{-\infty}^{\infty}|f(t)|dt<\infty\)。

绝对可积是个强限制。阶跃信号、增长的指数、许多实际系统的冲激响应都不满足它，傅里叶变换对它们失效或需要引入广义函数。

\hypertarget{ux62c9ux666eux62c9ux65afux53d8ux6362ux628aux57faux653eux5bbdux7528ux8870ux51cfux56e0ux5b50ux538bux5236ux53d1ux6563}{%
\subsubsection{3.2
拉普拉斯变换：把基放宽，用衰减因子压制发散}\label{ux62c9ux666eux62c9ux65afux53d8ux6362ux628aux57faux653eux5bbdux7528ux8870ux51cfux56e0ux5b50ux538bux5236ux53d1ux6563}}

拉普拉斯变换把基从 \(e^{j\omega t}\) 放宽为
\(e^{st}=e^{\sigma t}e^{j\omega t}\)。写成被积函数的形式：

\[
F(s)=\int_{-\infty}^{\infty}\left[f(t)e^{-\sigma t}\right]e^{-j\omega t}dt.
\]

这表示：先把 \(f(t)\) 乘以衰减因子
\(e^{-\sigma t}\)，再对乘积做傅里叶变换。对右边信号，只要 \(\sigma\)
取得足够大，\(f(t)e^{-\sigma t}\) 就能被压成绝对可积------这正是 1.3 中
\(\sigma>a\) 的收敛条件。对左边信号，结论方向相反：衰减因子要在
\(t\to-\infty\) 的方向上压住信号，需要 \(\sigma\)
取得足够小，收敛条件形如
\(\sigma<a\)。对双边信号，两个方向都要压住，可行的 \(\sigma\)
必须同时满足下界与上界，是一个带状区间
\(a<\sigma<b\)。三种情形合起来，正好解释了 2.4
中收敛域为什么有右侧、左侧、带状三种形状。

拉普拉斯变换用''先衰减再投影''的办法，把一大批傅里叶变换处理不了的信号纳入变换框架。付出的代价是变换依赖
\(\sigma\) 的取值，且并非所有 \(\sigma\) 都可行，可行的 \(\sigma\)
范围就是收敛域。收敛域因此不是外加的繁文缛节，而是''衰减因子必须足够强''这一要求的直接记录。

\hypertarget{ux5085ux91ccux53f6ux662fux62c9ux666eux62c9ux65afux5728ux865aux8f74ux4e0aux7684ux622aux9762}{%
\subsubsection{3.3
傅里叶是拉普拉斯在虚轴上的截面}\label{ux5085ux91ccux53f6ux662fux62c9ux666eux62c9ux65afux5728ux865aux8f74ux4e0aux7684ux622aux9762}}

在虚轴上 \(s=j\omega\)，即
\(\sigma=0\)，拉普拉斯变换退化为傅里叶变换：\(F(s)\big|_{s=j\omega}=F(j\omega)\)。但这个代入有一个前提：虚轴必须落在收敛域内，否则
\(s=j\omega\)
处的积分根本不存在。由此得到一个贯穿控制论的等价关系。先交代前提：讨论对象是
LTI 系统，``稳定''指 BIBO 稳定，即有界输入产生有界输出；对 LTI
系统，BIBO 稳定等价于冲激响应绝对可积。在此前提下：

\[
\text{系统稳定}\iff\text{收敛域包含虚轴}\iff\text{频率响应存在}.
\]

第一条等价：BIBO 稳定等价于冲激响应绝对可积，即 \(h(t)e^{-0\cdot t}\)
可积，即 \(\sigma=0\)
是可行的衰减因子，即虚轴在收敛域内。第二条等价：频率响应 \(H(j\omega)\)
就是 \(H(s)\)
在虚轴上的取值，虚轴在收敛域内它才有定义。这就是''只有稳定系统谈频率响应才有意义''的严格理由。2.5
的情形二中，带状收敛域 \(-1<\mathrm{Re}(s)<2\)
包含虚轴，所以那个双边系统虽然非因果，却是稳定的，频率响应存在；情形一与情形三的收敛域都不含虚轴，不稳定，\(H(j\omega)\)
无意义。

\hypertarget{ux66f4ux6df1ux4e00ux5c42ux7684ux7edfux4e00est-ux662fux7279ux5f81ux51fdux6570}{%
\subsubsection{\texorpdfstring{3.4 更深一层的统一：\(e^{st}\)
是特征函数}{3.4 更深一层的统一：e\^{}\{st\} 是特征函数}}\label{ux66f4ux6df1ux4e00ux5c42ux7684ux7edfux4e00est-ux662fux7279ux5f81ux51fdux6570}}

拉普拉斯变换之所以能把微分方程变成代数方程，根本原因是 \(e^{st}\)
是微分算子的特征函数：

\[
\frac{d}{dt}e^{st}=s\,e^{st}.
\]

微分算子作用在 \(e^{st}\) 上，结果仍是 \(e^{st}\) 本身，只乘了一个数
\(s\)。这个数 \(s\) 就是特征值。推而广之，\(e^{st}\)
是一切线性时不变（LTI）系统的特征函数：设 LTI 系统的冲激响应为
\(h(t)\)，输入 \(e^{st}\)，输出为卷积

\[
y(t)=\int_{-\infty}^{\infty}h(\tau)e^{s(t-\tau)}d\tau=e^{st}\int_{-\infty}^{\infty}h(\tau)e^{-s\tau}d\tau=H(s)e^{st}.
\]

输入是 \(e^{st}\)，输出仍是 \(e^{st}\)，只乘了复数 \(H(s)\)。\(H(s)\)
正是冲激响应的拉普拉斯变换，也就是系统对特征函数 \(e^{st}\) 的特征值。

这解释了变换方法的全部威力。微分方程之所以难解，是因为求导把不同函数混在一起；但
\(e^{st}\) 在求导下保持不变（只差一个倍数），把信号分解到 \(e^{st}\)
这组基上，微分算子在每个基方向上的作用就退化为乘以
\(s\)，微分方程变成关于 \(s\)
的代数方程。同理，卷积是输入与冲激响应的混合，但在 \(e^{st}\)
基下，卷积退化为特征值相乘：\(Y(s)=H(s)U(s)\)。这正是线性代数中特征分解思想在函数空间中的对应：找到算子的特征向量，算子的作用就退化为标量乘法。

\hypertarget{ux56dbux4f20ux9012ux51fdux6570}{%
\subsection{四、传递函数}\label{ux56dbux4f20ux9012ux51fdux6570}}

\hypertarget{ux5b9aux4e49ux4e0eux57faux672cux6027ux8d28}{%
\subsubsection{4.1
定义与基本性质}\label{ux5b9aux4e49ux4e0eux57faux672cux6027ux8d28}}

对零初始条件的 LTI 系统，定义传递函数

\[
G(s)=\frac{Y(s)}{U(s)},
\]

其中 \(U(s)\)、\(Y(s)\)
分别为输入、输出的拉普拉斯变换。零初始条件是定义的一部分：它把初始储能的影响排除在外，使
\(G(s)\) 只反映系统本身。

传递函数有四条基本性质。第一，它是系统的固有属性，与输入信号的形式、大小无关，也与初始条件无关。第二，它是冲激响应
\(g(t)\) 的拉普拉斯变换：令 \(u(t)=\delta(t)\)，则
\(U(s)=1\)，\(Y(s)=G(s)\)，故 \(y(t)=g(t)\)。第三，时域卷积对应 \(s\)
域乘法：\(y(t)=g(t)*u(t)\) 等价于 \(Y(s)=G(s)U(s)\)。第四，由
3.4，\(G(s)\) 是系统对特征函数 \(e^{st}\) 的特征值。

\hypertarget{ux80fdux89e3ux65b9ux7a0bux4e3aux4ec0ux4e48ux8fd8ux8981ux4f20ux9012ux51fdux6570}{%
\subsubsection{4.2
能解方程，为什么还要传递函数}\label{ux80fdux89e3ux65b9ux7a0bux4e3aux4ec0ux4e48ux8fd8ux8981ux4f20ux9012ux51fdux6570}}

微分方程模型已经完整描述了系统，配上方程解法原则上可以求出任何响应。传递函数的必要性体现在五个方面：

\begin{enumerate}
\def\labelenumi{\arabic{enumi}.}
\tightlist
\item
  方框图代数。系统互联后，求总模型不再是联立微分方程组，而是对传递函数做四则运算（见第五节）。这是工程上最实质的收益。
\item
  不解方程直接读特性。在控制论的因果默认下，\(G(s)\)
  的零极点分布直接给出稳定性与动态特性，无需解出
  \(y(t)\)。脱离这个默认则不然：同一组极点配上不同的收敛域，稳定性结论可以完全不同，2.5
  的三种情形正是例证。
\item
  频率响应。令 \(s=j\omega\) 得
  \(G(j\omega)\)，其模与相角给出正弦稳态响应的幅值增益与相移，是波特图与频域设计的基础。
\item
  控制器设计的对象。控制问题的标准提法是：选定控制器
  \(C(s)\)，使闭环极点落在指定位置。设计变量是传递函数，不是微分方程的解。
\item
  卷积变乘法。复杂输入下的响应计算由卷积积分变为代数乘法再加一次反变换。
\end{enumerate}

关于第 3 点还需说明 \(G(j\omega)\) 为什么给出正弦稳态响应。由 3.4，输入
\(e^{st}\) 产生输出 \(G(s)e^{st}\)。取 \(s=j\omega\)，输入
\(e^{j\omega t}\) 产生输出 \(G(j\omega)e^{j\omega t}\)。正弦输入
\(\sin\omega t\) 是 \(e^{j\omega t}\) 与 \(e^{-j\omega t}\)
的线性组合，由线性性，稳态输出是同频率正弦的线性组合，幅值被放大
\(|G(j\omega)|\) 倍，相位移动
\(\angle G(j\omega)\)。前提是系统稳定（收敛域含虚轴），\(G(j\omega)\)
才有定义，且暂态分量衰减后剩下的才是稳态。

\hypertarget{ux96f6ux70b9ux4e0eux6781ux70b9}{%
\subsubsection{4.3 零点与极点}\label{ux96f6ux70b9ux4e0eux6781ux70b9}}

把传递函数写成多项式之比：

\[
G(s)=\frac{N(s)}{D(s)}.
\]

\textbf{零点}是 \(N(s)=0\) 的根。零点 \(s=z\) 的准确含义是：系统对
\(e^{zt}\) 型输入的稳态传输为零。由 3.4，输入 \(e^{zt}\) 产生的输出为
\(G(z)e^{zt}=0\)，即这个指数成分被系统完全阻断。特别地，只有位于虚轴上的零点
\(z=j\omega_0\) 才对应阻断某个正弦频率成分：\(G(j\omega_0)=0\)
意味着频率为 \(\omega_0\)
的正弦输入的稳态输出为零；不在虚轴上的零点阻断的是 \(e^{zt}\)
型的指数或阻尼振荡成分，而不是某个正弦频率。零点还决定各自然模态在输出中的配比------不同的零点分布会让同样一组模态以不同的权重组合，从而改变响应的具体形状。

\textbf{极点}是 \(D(s)=0\) 的根。\(s\) 趋于极点时
\(|G(s)|\to\infty\)；极点本身处 \(G(s)\) 没有定义，不能说''\(G(s)\)
在极点等于无穷''。极点决定系统的自然模态：每个极点 \(p\)
对应响应中的一个成分
\(e^{pt}\)。极点还决定稳定性与动态特性：极点全部在左半平面，所有模态衰减，系统稳定；共轭复极点对应
\(e^{\sigma t}\sin(\omega t+\phi)\)
型模态，响应振荡；极点越靠近虚轴，对应模态衰减越慢，响应越慢。

一个常见的错误是把极点当成''\(G(s)=0\) 的解''。必须纠正：\(G(s)=0\)
的解是零点；极点是 \(G(s)\) 趋于无穷的解，即分母 \(D(s)=0\)
的根。零点是分子为零之处，极点是分母为零之处，二者对系统的作用完全不同。

\hypertarget{ux4e94ux65b9ux6846ux56feux4ee3ux6570}{%
\subsection{五、方框图代数}\label{ux4e94ux65b9ux6846ux56feux4ee3ux6570}}

\hypertarget{ux57faux672cux5143ux4ef6ux4e0eux4e09ux6761ux89c4ux5219}{%
\subsubsection{5.1
基本元件与三条规则}\label{ux57faux672cux5143ux4ef6ux4e0eux4e09ux6761ux89c4ux5219}}

方框图有两种基本元件。方框：输入 \(U(s)\) 经方框 \(G(s)\) 后输出
\(G(s)U(s)\)。相加点：输出为各输入的代数和，符号在相加点旁标注。

由这两种元件可以建立三条等效规则。

\textbf{串联。} 两个方框 \(G_1\)、\(G_2\) 串联，中间信号为
\(G_1U\)，输出为 \(G_2G_1U\)，等效传递函数为 \(G_1G_2\)。

\textbf{并联。} 两路 \(G_1\)、\(G_2\)
接受同一输入，输出在相加点相加：\(Y=G_1U+G_2U=(G_1+G_2)U\)，等效传递函数为
\(G_1+G_2\)。

\textbf{负反馈。} 前向通道 \(G(s)\)，反馈通道 \(H(s)\)，输出 \(Y\) 经
\(H\) 反馈回相加点，与输入 \(U\) 相减。设相加点输出（误差信号）为
\(E\)，则

\[
E=U-HY,\qquad Y=GE.
\]

把第一式代入第二式：

\[
Y=G(U-HY)=GU-GHY,
\]

\[
Y+GHY=GU,\qquad Y(1+GH)=GU,
\]

\[
\frac{Y}{U}=\frac{G}{1+GH}.
\]

若相加点处取加号（正反馈），同样推导得
\(\dfrac{Y}{U}=\dfrac{G}{1-GH}\)，分母为 \(1-GH\)。

复杂方框图的化简方法是：找到最内层能套用三条规则之一的局部结构，等效为一个方框，逐层由内向外推进，直到全图化为一个传递函数。

\hypertarget{ux5b8cux6574ux95edux73afux793aux4f8b}{%
\subsubsection{5.2
完整闭环示例}\label{ux5b8cux6574ux95edux73afux793aux4f8b}}

设控制器 \(C(s)=K\)（比例控制），被控对象

\[
G(s)=\frac{1}{s(s+1)},
\]

单位负反馈（\(H(s)=1\)）。开环传递函数为

\[
C(s)G(s)=\frac{K}{s(s+1)}.
\]

代入负反馈公式，闭环传递函数

\[
T(s)=\frac{C(s)G(s)}{1+C(s)G(s)}=\frac{\dfrac{K}{s(s+1)}}{1+\dfrac{K}{s(s+1)}}.
\]

分子分母同乘 \(s(s+1)\)：

\[
T(s)=\frac{K}{s(s+1)+K}=\frac{K}{s^{2}+s+K}.
\]

现在比较开环与闭环的极点。开环 \(C(s)G(s)\) 的极点是 \(s=0\) 与
\(s=-1\)，即对象本身的极点。闭环 \(T(s)\) 的极点是特征方程

\[
s^{2}+s+K=0
\]

的根：

\[
s=\frac{-1\pm\sqrt{1-4K}}{2}.
\]

闭环极点随 \(K\) 移动，分三种情形。

取 \(K<0.25\)：判别式 \(1-4K>0\)，极点为两个不相等的负实根

\[
s_1=\frac{-1+\sqrt{1-4K}}{2},\qquad s_2=\frac{-1-\sqrt{1-4K}}{2},
\]

系统过阻尼。两个根一快一慢：\(s_2\) 更靠左，对应模态 \(e^{s_2 t}\)
衰减得快；\(s_1\) 更靠近虚轴，对应模态 \(e^{s_1 t}\)
衰减得慢。响应是两个指数衰减的叠加，初始阶段受快模态影响，时间一长快模态消失，剩下的过程几乎完全由慢模态决定------慢极点
\(s_1\) 主导响应，调节时间由它说了算。\(K\) 越小，\(s_1\)
越靠近原点，响应越迟缓。

取 \(K=0.25\)：判别式 \(1-4K=0\)，得重根
\(s=-0.5\)，系统临界阻尼，在不过冲的前提下调节时间最短。

取 \(K>0.25\)：判别式为负，极点为共轭复数

\[
s=-\frac{1}{2}\pm j\frac{\sqrt{4K-1}}{2},
\]

实部固定为 \(-\frac{1}{2}\)，虚部随 \(K\)
增大，响应振荡且振荡频率升高。\(K\) 越大，振荡越明显。

把三种情形连起来，可以看出 \(K\) 从 \(0\) 增大时闭环极点在 \(s\)
平面上的移动过程。\(K=0\) 时极点为 \(0\) 与
\(-1\)，即开环极点本身。\(K\)
从零开始增大，两个极点沿实轴相向移动，响应保持过阻尼但慢极点左移、响应加快；\(K=0.25\)
时两极点在 \(s=-0.5\) 处汇合为重根，到达临界阻尼；\(K\)
继续增大，极点离开实轴，分裂为实部固定为 \(-\frac{1}{2}\)
的共轭复根，沿竖直线 \(\mathrm{Re}(s)=-\frac{1}{2}\)
向上下两个方向分离，系统进入欠阻尼，振荡越来越强。实轴上的相向移动对应过阻尼到临界阻尼的过渡，离开实轴则标志振荡的出现。

这个例子说明了三件事。第一，闭环极点不等于开环极点：开环极点是 \(0\) 与
\(-1\)，闭环极点是 \(s^2+s+K=0\)
的根，反馈改变了系统的模态。第二，调参的实质是移动闭环极点：改变 \(K\)
就是改变特征方程的系数，从而改变极点位置，进而改变响应的阻尼与振荡特性。第三，让闭环极点随参数
\(K\) 从 \(0\) 到 \(\infty\) 移动并在 \(s\)
平面上描出轨迹，正是根轨迹法的起点。

\hypertarget{ux516dux5e94ux7528ux793aux4f8bux4e24ux8f6eux5e73ux8861ux8f66ux7684ux7cfbux7edfux63a7ux5236ux6846ux56fe}{%
\subsection{六、应用示例：两轮平衡车的系统控制框图}\label{ux516dux5e94ux7528ux793aux4f8bux4e24ux8f6eux5e73ux8861ux8f66ux7684ux7cfbux7edfux63a7ux5236ux6846ux56fe}}

前面第五节的方框图代数可以直接用于两轮平衡车。网上已有比较成熟的画法：一类仿真框图以期望姿态角
\(0^\circ\) 为输入，前向通道依次是 PID
控制器、电机模型、车体模型，输出为姿态角 \(\theta\)、位移 \(x\) 和角速度
\(\dot\theta\)；实际硬件框图则通常把主控、IMU、编码器、电机驱动和左右电机构成闭环。\footnote{詹友基,
  应廉泽, 戴福全. 两轮自平衡车控制系统的研究{[}J{]}. 福建工程学院学报,
  2018, 16(6): 511--515. DOI: 10.3969/j.issn.1672-4348.2018.06.001.}\footnote{Design
  and implementation of two-wheeled self-balancing mobile robot control
  system based on STM32-MAT and Android{[}EB/OL{]}. 2025.
  https://www.extrica.com/article/24881.}
按题目要求，这里不展开具体传递函数，只明确各个环节及其输入、输出量。

\begin{figure}[H]
\centering
\resizebox{\textwidth}{!}{%
\begin{tikzpicture}[
    >=Latex,
    block/.style={draw, rounded corners=2pt, minimum width=1.75cm, minimum height=1.05cm, align=center, font=\small, fill=blue!4},
    sensor/.style={draw, rounded corners=2pt, minimum width=2.15cm, minimum height=1.05cm, align=center, font=\small, fill=green!6},
    sumnode/.style={circle, draw, minimum size=0.68cm, inner sep=0pt, font=\small},
    signal/.style={-{Latex[length=2.2mm]}, thick},
    lab/.style={font=\scriptsize, align=center}
]
\node[block] (ref) {期望姿态角\\$\theta_{\mathrm{ref}}=0^\circ$};
\node[sumnode, right=0.85cm of ref] (sum) {$\Sigma$};
\node[block, right=0.95cm of sum] (pid) {姿态 PID\\控制器};
\node[block, right=0.85cm of pid] (driver) {电机驱动\\PWM/电压};
\node[block, right=0.85cm of driver] (motor) {左右直流\\电机};
\node[block, right=0.85cm of motor, minimum width=2.15cm] (plant) {两轮平衡车\\车体（倒立摆）};
\node[right=0.9cm of plant, align=left, font=\small, text width=2.55cm] (out) {\textbf{输出量}\\$\theta(t)$：姿态角\\$x(t)$：位移\\$\dot\theta(t)$：角速度};

\draw[signal] (ref) -- (sum) node[midway, above, lab] {参考输入};
\draw[signal] (sum) -- (pid) node[midway, above, lab] {$e_\theta$};
\draw[signal] (pid) -- (driver) node[midway, above, lab] {$u(t)$};
\draw[signal] (driver) -- (motor);
\draw[signal] (motor) -- (plant) node[midway, above, lab] {转矩/转速};
\draw[signal] (plant) -- (out);

\node[lab] at ($(sum.west)+(-0.18,0.18)$) {$+$};
\node[lab] at ($(sum.south)+(-0.18,-0.18)$) {$-$};

\node[above=0.7cm of plant, font=\small] (dist) {外部扰动 $d(t)$};
\draw[signal] (dist) -- (plant.north);

\node[sensor, below=1.55cm of motor] (measure) {姿态/速度测量\\IMU + 编码器};
\node[sensor, below=1.55cm of pid, minimum width=2.35cm] (filter) {滤波与姿态解算\\互补/卡尔曼滤波};
\draw[signal] (plant.south) |- (measure.east);
\draw[signal] (measure) -- (filter) node[midway, above, lab] {原始测量};
\draw[signal] (filter.west) -| (sum.south) node[pos=0.28, below, lab] {$\hat\theta,\hat{\dot\theta},\hat v$};
\end{tikzpicture}}
\caption{两轮平衡车系统控制框图（按典型 PID 闭环结构重画）}
\label{fig:balance-car-control}
\end{figure}

图 \ref{fig:balance-car-control} 中各环节的含义如下：

\begin{enumerate}
\def\labelenumi{\arabic{enumi}.}
\tightlist
\item
  \textbf{输入量}：期望姿态角
  \(\theta_{\mathrm{ref}}\)，直立平衡时通常取
  \(0^\circ\)；若考虑前进、后退，还可另给定期望速度或期望位移。
\item
  \textbf{控制器}：常用 PID 或 PD 控制。它根据姿态误差
  \(e_\theta=\theta_{\mathrm{ref}}-\hat\theta\) 产生控制量
  \(u(t)\)，在嵌入式实现中通常表现为 PWM 占空比或等效电机电压。
\item
  \textbf{执行机构}：电机驱动器与左右直流电机，把控制量转换为车轮的转矩和转速。
\item
  \textbf{被控对象}：两轮平衡车车体。小倾角下可近似看作``可移动底座上的倒立摆''，其天然开环特性是不稳定的。
\item
  \textbf{输出量}：车体姿态角 \(\theta(t)\)、角速度
  \(\dot\theta(t)\)、车轮位移 \(x(t)\) 或速度 \(v(t)\)。
\item
  \textbf{反馈通道}：IMU（加速度计与陀螺仪）测量姿态，编码器测量轮速；测量值经滤波与姿态解算后得到
  \(\hat\theta\)、\(\hat{\dot\theta}\)、\(\hat v\)，再送回相加点形成负反馈。
\end{enumerate}

从控制结构上看，它与第五节的单位负反馈完全一致：相加点产生误差，控制器根据误差产生控制量，车体输出再经传感器反馈。实际平衡车还会在姿态环外叠加速度环和方向环，但本章只保留最能说明控制结构的总体框图。

\hypertarget{ux4e03ux5982ux4f55ux68c0ux9a8cux81eaux5df1ux662fux5426ux771fux61c2}{%
\subsection{七、如何检验自己是否真懂}\label{ux4e03ux5982ux4f55ux68c0ux9a8cux81eaux5df1ux662fux5426ux771fux61c2}}

控制论训练的价值不在执行计算，而在三处：建模------把物理对象写成微分方程或传递函数；解读------从零极点分布读出稳定性与动态特性；设计------选择控制器使闭环极点落在指定位置。这三件事都要求理解变换的含义，而不是会算变换。

如果只是能够流畅地复述本文的每一句话，会给人''已经懂了''的错觉。理解与否的检验标准是闭卷产出------不翻书、不查表，独立完成推导、做题与讲解。以下清单供闭卷自测：

\begin{enumerate}
\def\labelenumi{\arabic{enumi}.}
\tightlist
\item
  从积分定义出发，推出 \(u(t)\) 与 \(e^{at}u(t)\)
  的拉普拉斯变换，并写出各自的收敛条件。
\item
  独立推导负反馈公式 \(\dfrac{Y}{U}=\dfrac{G}{1+GH}\)。
\item
  对 \(H(s)=\dfrac{1}{(s+1)(s-2)}\)，独立写出三种收敛域对应的
  \(h(t)\)，并指出哪种收敛域对应稳定系统。
\item
  解释为什么只有稳定系统的频率响应才有意义（从收敛域含虚轴说起）。
\item
  说明为什么 \(G(j\omega)\) 给出正弦稳态响应（从 \(e^{st}\) 是 LTI
  系统的特征函数说起）。
\end{enumerate}

五条中任何一条做不到，回到对应章节重读推导，再闭卷做一次，直到全部通过。

\end{document}
